%!TEX program = xelatex
%!TEX encoding = ISO-8859-1
\documentclass[12pt]{article}

\usepackage{amssymb}
\usepackage{amsmath,amsfonts,amsthm,amstext,mathabx}
\usepackage[brazil]{babel}
%\usepackage[latin1]{inputenc}
\usepackage{graphicx}
\graphicspath{{/home/jfreitas/Dropbox/imagens-latex/}{/ArquivosUbuntu/Dropbox/imagens-latex/}{D:/Dropbox/imagens-latex/}}
\usepackage{enumitem}
\usepackage{multicol}
\usepackage[all]{xy}

\setlength{\topmargin}{-1.0in}
\setlength{\oddsidemargin}{0in}
\setlength{\textheight}{10.1in}
\setlength{\textwidth}{6.5in}
\setlength{\baselineskip}{12mm}

\newcounter{exercicios}
\renewcommand{\theexercicios}{\arabic{exercicios}}

\newenvironment{questao}[1]{%
\refstepcounter{exercicios}%
\label{#1}%
\noindent\textbf{Exerc{\'\i}cio {\theexercicios}:}%
}%
{}%

% \newcounter{exercicios}
% \setcounter{exercicios}{0}
% \newcommand{\questao}{
% \addtocounter{exercicios}{1}
% \noindent{\bf Exerc{\'\i}cio \arabic{exercicios}: }}

\newcommand{\equi}{\Leftrightarrow}
\newcommand{\bic}{\leftrightarrow}
\newcommand{\cond}{\rightarrow}
\newcommand{\impl}{\Rightarrow}
\newcommand{\nao}{\sim}
\newcommand{\sub}{\subseteq}
\newcommand{\e}{\ \wedge\ }
\newcommand{\ou}{\ \vee\ }
\newcommand{\vaz}{\emptyset}
\newcommand{\nsub}{\nsubseteq}
\renewcommand{\sin}{{\rm sen\,}}

\newcommand{\n}{\mathbb{N}}
\newcommand{\z}{\mathbb{Z}}
\newcommand{\q}{\mathbb{Q}}
\newcommand{\im}{{\rm Im\,}}
\newcommand{\aut}{{\rm Aut\,}}
\newcommand{\cp}[1]{\mathbb{#1}}
\newcommand{\real}{\mathbb{R}}
\newcommand{\complex}{\mathbb{C}}
\newcommand{\lap}[1]{\mathcal{L}\left\{#1\right\}}
\newcommand{\lapi}[1]{\mathcal{L}^{-1}\left\{#1\right\}}
\newcommand{\se}[1]{\displaystyle\sum_{n = 1}^\infty{#1}}
\newcommand{\dlim}[2]{\displaystyle\lim_{#1\rightarrow #2}}
\newcommand{\slim}{\displaystyle\lim_{n \rightarrow \infty}}
\newcommand{\seq}[1]{\{{#1_n\}}}
\newcommand{\seg}[1]{\displaystyle\sum_{n = 1}^\infty{#1_n}}
\newcommand{\sei}[2]{\displaystyle\sum_{#1}^\infty{#2}}
\newcommand{\sepc}[3]{\displaystyle\sum_{#1}^\infty{#2(x - #3)^n}}
\newcommand{\imp}[3]{\displaystyle\int_{#1}^{+\infty}{#3}{d #2}}
\newcommand{\dint}[4]{\displaystyle\int_{#1}^{#2}{#4}{d#3}}
\newcommand{\inti}[2]{\displaystyle\int{#1}{d#2}}
\newcommand{\norma}[1]{\left\lVert#1\right\rVert}
\newcommand{\flim}[1]{\displaystyle\lim_{#1\rightarrow \infty}}
\renewcommand{\sin}{{\rm sen\,}}
\renewcommand{\tan}{{\rm tg\,}}
\renewcommand{\csc}{{\rm cossec\,}}
\renewcommand{\cot}{{\rm cotg\,}}
\renewcommand{\sinh}{{\rm senh\,}}
\newcommand\T{\rule{0pt}{2.6ex}} 
\newcommand{\vesp}{\vspace{0.2cm}}
\newcommand{\subne}{\subsetneqq}


\newcommand{\compcent}[1]{\vcenter{\hbox{$#1\circ$}}}
\newcommand{\comp}{\mathbin{\mathchoice
{\compcent\scriptstyle}{\compcent\scriptstyle}
{\compcent\scriptscriptstyle}{\compcent\scriptscriptstyle}}}

\begin{document}

\pagestyle{empty}

\begin{figure}[h]
        \begin{minipage}[c]{1.7cm}
        \includegraphics[width=1.7cm]{unb.pdf}
        \end{minipage}%
        \hspace{0pt}
        \begin{minipage}[c]{4in}
          {Universidade de Bras{\'\i}lia} \\
          {Departamento de Matem{\'a}tica}
\end{minipage}
\end{figure}
\vspace{-1cm}\hrule


\begin{center}
{\Large\bf {\'A}lgebra 1 - Turma C -- 1$^{o}$/2018} \\ \vspace{9pt} {\large\bf
  $3^{\underline{o}}$ Teste - Resolu\c{c}\~ao}\\
\vspace{9pt} Prof. Jos{\'e} Ant{\^o}nio O. Freitas
\end{center}
\hrule

\vspace{.6cm}

\questao Para mostrar que $R$ é relação de equivalência sobre $A$ precisamos mostrar que
\begin{enumerate}
	\item Para todo $x \in A$, $xRx$.
	\item Se $xRy$, então $yRx$.
	\item Se $xRy$ e $yRz$, então $XRz$.
\end{enumerate}

Como $A = \n \times \n^*$, então os elementos que aparecem na definição de relação de equivalência serão pares ordenados. Nesse caso particular precisamos verificar que:
\begin{enumerate}
	\item Para todo $x = (a,b) \in A$, $(a,b)R(a,b)$.
	\item Para $x = (a,b)$ e $y = (c,d)$, se $(a,b)R(c,d)$, então $(c,d)R(a,b)$.
	\item Para $x = (a,b)$, $y = (c,d)$ e $z = (e,f)$ se $(a,b)R(c,d)$ e $(c,d)R(e,f)$, então $(a,b)R(e,f)$.
\end{enumerate}

De fato,
\begin{enumerate}
	\item Para $(a,b) \in A$ $a + b = a + b$, logo $(a,b)R(a,b)$.
	\item Se $(a,b)R(c,d)$ então, por definição $a + b = c + d$. Ou seja $c + d = a + b$ e com isso $(c,d)R(a,b)$.
	\item se $(a,b)R(c,d)$ e $(c,d)R(e,f)$ então, por definição $a + b = c + d$  e $c + d = e + f$. Logo $a + b = e + f$, isto é, $(a,b)R(e,f)$.
\end{enumerate}

Portanto, $R$ é uma relação de equivalência sobre $A$ como queríamos.

\[
	(a,b)R(c,d) \mbox{ se, e só se, } a + b = c + d.
\]
Mostre que $R$ é uma relação de equivalência sobre $A$.

\vspace{.5cm}

\questao Considere a seguinte relação sobre $\z$
\[
	xRy \mbox{ quando } x^2 - y^2 = 4k, \mbox{para algum} k \in \z.
\]

\begin{enumerate}[label={\alph*})]
	\item Mostre que $R$ é uma rela{\c c}{\~a}o de equival{\^e}ncia sobre $\z$.
	\item Encontre $\overline{0}$, $\overline{1}$ e $\overline{2}$.
\end{enumerate}

\noindent\textbf{Solu\c{c}\~ao:} De fato,
	\begin{enumerate}[label={\alph*})]
		\item Primeiro observe que
		\[
			R = \{((a,b);(c,d)) \in A \times A \mid ad = bc\}
		\]
		Assim um elemento $x \in A = \z\times\z^*$ é da forma $x = (a,b)$.
		\begin{enumerate}[label={\roman*})]
			\item Para $(a,b) \in A$, $ab = ab$ logo $(a,b)R(a,b)$.
			\item Se $(a,b)R(c,d)$, então $ad = bc$. Assim $bc = ad$, isto é, $(c,d)R(a,b)$.
			\item Se $(a,b)R(c,d)$ e $(c,d)R(e,f)$ então
			\begin{align}
				ad &= bc\label{primeira_equacao}\\
				cf &= de.\label{segunda_equacao}
			\end{align}
			Multiplicando a equação \eqref{primeira_equacao} por $f$ obtemos
			\begin{align}\label{terceira_equacao}
				adf = bcf.
			\end{align}
			Substituindo o valor de $cf$ da equação \eqref{segunda_equacao} na equação \eqref{terceira_equacao} obtemos
			\begin{align*}
				adf &= bde\\
				adf - bde &= 0\\
				d(af - be) &= 0.
			\end{align*}
			Como $d \in \z^*$, $d \ne 0$, logo $af - be = 0$. Assim $af = be$ e com isso $(a,b)R(e,f)$.
		\end{enumerate}

		Portanto $R$ é uma relação de equivalência sobre $A$.

		\item A classe de equivalência de $\overline{(2,1)}$ é dada por:
		\begin{align*}
			\overline{(2,1)} &= \{(x,y) \in A \mid (x,y)R(2,1)\} = \{(x,y) \in A \mid x = 2y\} \\
			&= \{\cdots; (2,1); (4,2); (-2,-1);(-4,-2);(6,3);\cdots\}.
		\end{align*}
	\end{enumerate}
\end{document}